\begin{pastexam}[label={px:ccu-115-swa-06}]{中正資工所甲組 115 學年度 第 16 題}
（複選；全對才給分）

(2\%) (\qquad) Suppose we are debugging a Quicksort implementation that is intended to sort an array
in ascending order. After the first partition step has been completed, the contents of the array
are:
\begin{center}5, 2, 8, 11, 20, 18, 14, 12\end{center}
Which of the following statements are \textbf{incorrect} about the Quicksort?
\begin{enumerate}[label=(\alph*)]
\item Quicksort is a stable sorting algorithm because elements equal to the pivot always preserve
their relative order during partitioning.
\item If the pivot is always chosen as the smallest or largest element, Quicksort still produces a
correctly sorted array, even though its running time becomes $O(n^2)$.
\item The pivot could have been 11, but could not have been 14.
\item Quicksort always runs in $O(n\log n)$ time because each partition step divides the array into
two nearly equal parts.
\end{enumerate}
\end{pastexam}

\begin{pastsoln}
答案：(a)(d)。 題目問的是「錯誤」的敘述。

\begin{itemize}
\item (a) 是錯的敘述，選。 一般的分割會把元素跨過很遠交換，等值元素的相對順序可能被打亂，
快速排序不是穩定的（資結書第~14.1、14.5~節）。 例如 Lomuto 分割（資結書虛擬碼~14.8）排
$\langle 2_a, 2_b, 1\rangle$：基準是最後的 $1$，兩個 $2$ 都不小於它，最後基準和 $2_a$ 交換，得到
$\langle 1, 2_b, 2_a\rangle$，兩個 $2$ 的順序反了。
\item (b) 是對的敘述，不選。 基準怎麼選只影響每次切得平不平均，不影響正確性：每次分割後基準
都在最終位置，左右兩段遞迴排好，整個就排好了。 每次都選到最大或最小，切成 $0$ 和 $m - 1$，
$T(n) = T(n - 1) + \Theta(n) = \Theta(n^2)$（資結書第~14.5~節）。
\item (c) 是對的敘述，不選。 分割完後，基準左邊都不大於它、右邊都不小於它。 逐一檢查：
\begin{itemize}
\item $11$（第 $4$ 個）：左邊 $5, 2, 8$ 都 $< 11$，右邊 $20, 18, 14, 12$ 都 $> 11$，可以是基準。
\item $14$（第 $7$ 個）：左邊有 $20, 18 > 14$，不可能是基準。
\end{itemize}
（順帶一提，$8$ 也滿足條件，也可能是基準；但 (c) 只說 $11$ 可能、$14$ 不可能，兩句都對。）
\item (d) 是錯的敘述，選。 「每次都切成差不多兩半」不是保證：(b) 就是反例，最壞情況
$\Theta(n^2)$。 $O(n \log n)$ 是平均情況（資結書定理~14.3）或隨機選基準時的期望時間。
\end{itemize}
\end{pastsoln}
